\mxcorrections{3}\label{C03-CORR}
\textbf{Usage professeur.} Les réponses qui suivent concernent les tâches non résolues dans le cours. Faire d’abord expliciter la démarche ; fournir une aide graduée avant la solution. Les durées sont des estimations à étalonner en classe.
\subsection*{Diagnostic et remédiation}
\begin{enumerate}
\item $\pt{3;-4}$. \textit{Action :} Soustraire arrivée moins départ sur deux dessins.
\item $\sqrt{3^2+(-4)^2}=5$. \textit{Action :} Revoir Pythagore et distinguer norme et coordonnées.
\item $\vect{AC}$. \textit{Action :} Tracer le déplacement composé.
\item $a^2+2ab+b^2$. \textit{Action :} Identifier le terme croisé.
\item $\dfrac{3}{5}$. \textit{Action :} Nommer les côtés par rapport à l’angle.
\item $\pt{1;2}$. \textit{Action :} Moyenner séparément les deux coordonnées.
\end{enumerate}
\subsection*{Activité de recherche}
Prendre des points sur le cercle de diamètre AB ; démonstration par le milieu I : produit $MI^2-9$. Les extrémités sont incluses car l’un des vecteurs est nul.
\subsection*{Applications immédiates}
\begin{enumerate}
\item $0$ puis $-10$. Le cosinus vaut respectivement 0 et $-1$.
\item $3t-12=0$, donc $t=4$.
\item Le carré vaut $4-2+9=11$ ; la norme vaut $\sqrt{11}$.
\item $32=16+16-32\cos A$, donc $\cos A=0$ et $A=90^\circ$.
\item $MI^2-9=7$ : cercle de centre $\pt{3;0}$ et de rayon 4.
\end{enumerate}
\subsection*{Exercices : démarche et résultats}
\paragraph{Exercice 3.1 -- Calculs directs.} Produit $4-6=-2$ ; normes $\sqrt5$ et 5 ; somme $\pt{5;1}$ de carré 26. Contrôle : $5+2(-2)+25=26$.
\paragraph{Exercice 3.2 -- Un triangle rectangle ?.} $\vect{AB}=\pt{3;1}$, $\vect{AC}=\pt{-1;3}$ ; produit 0. $AB=AC=\sqrt{10}$, aire $\dfrac{AB\times AC}{2}=5$ unités d’aire.
\paragraph{Exercice 3.3 -- Longueur sans repère.} $\vect{AC}=\vect{AB}+\vect{AD}$ et $\vect{BD}=\vect{AD}-\vect{AB}$. Produit $\dfrac{4\times 3}{2}=6$ ; $AC^2=16+9+12=37$ et $BD^2=16+9-12=13$.
\paragraph{Exercice 3.4 -- Retrouver un angle.} $\cos \theta =-\dfrac{6}{12}=-\dfrac{1}{2}$ ; dans $\intff{0}{180^\circ}$, l’angle est $120^\circ$. Le domaine assure l’unicité.
\paragraph{Exercice 3.5 -- Un lieu variable.} $MI^2=k+\dfrac{AB^2}{4}$. Si le second membre est négatif : ensemble vide ; nul : le seul point I ; positif : cercle de rayon sa racine carrée. La discussion porte sur une distance au carré.
\paragraph{Exercice 3.6 -- Travail d’une force.} $W=\dfrac{50\times 4\times 1}{2}=100$ J. À $120^\circ$, $W=-100$ J. Le signe distingue l’aide au déplacement de l’opposition ; il ne décrit pas une distance négative.
\paragraph{Exercice 3.7 -- Contraposée.} La propriété directe « orthogonaux implique égalité des carrés » résulte du développement. Sa contraposée est exactement l’énoncé. On peut aussi écrire $2\vec u\cdot\vec v\ne0$.
\paragraph{Exercice 3.8 -- Reprise différée.} Al-Kashi : $\dfrac{9+16-25}{2}=0$. Réciproque de Pythagore : triangle rectangle en A, donc produit nul. Nommer les hypothèses de chaque méthode.
\paragraph{CQFD, exercice 9 p. 250, questions a et b.} a : $1\times3+2\times(-1)=1$. b : $(-1)\times12+3\times7=9$. Le produit scalaire est un nombre, pas un vecteur.
\paragraph{CQFD, exercice 11 p. 250, questions 1, 2 et 3.} $\vect{AB}=\pt{4;6}$, $\vect{AC}=\pt{-3;4}$, $I\pt{-1;-2}$ et $\vect{CI}=\pt{5;-1}$. D’où $\vect{AB}\cdot\vect{AC}=12$, $\vect{AB}\cdot\vect{CI}=14$. Chasles donne $\vect{AB}\cdot\vect{AI}=12+14=26$. Deux contrôles : $\vect{AI}=\pt{2;3}$ donne $8+18=26$ ; ou $\vect{AI}=\dfrac12\vect{AB}$ donne $\dfrac12(16+36)=26$.
\subsection*{QCM : réponses et raisons}
\textbf{QCM 1.}
\begin{enumerate}
\item A. $6-4=2$.
\item C. Racine de 25.
\item D. Critère scalaire.
\item B. $\dfrac{2\times 3}{2}=3$.
\item C. Développer le carré scalaire.
\end{enumerate}
\textbf{QCM 2.}
\begin{enumerate}
\item D. Cosinus négatif.
\item A. Coordonnées par projection.
\item C. $MI^2=\dfrac{AB^2}{4}$.
\item B. Bilinéarité : $2(-1)3$.
\item D. $BC^2=8$.
\end{enumerate}
\subsection*{Vrai ou faux}
\begin{enumerate}
\item Faux : $u\cdot(-u)=-\norm u^2<0$ pour $u\ne0$.
\item Vrai : son produit avec tout vecteur vaut 0.
\item Faux : $\pt{1;0}$ et $\pt{0;1}$ sont non nuls et orthogonaux.
\end{enumerate}
\paragraph{Différenciation.} Pour un blocage technique, proposer la première étape seulement (domaine, facteur ou univers de référence), puis retirer l’aide. Pour un élève autonome, demander une seconde méthode et le traitement d’un cas limite. À l’oral, autoriser une reformulation initiale dans la langue maîtrisée, puis exiger une phrase mathématique précise en français.
